\documentclass[11pt,a4paper]{article}
\usepackage[left=2cm,right=2cm,top=2.5cm,bottom=2.5cm]{geometry}
\usepackage{fontspec}
\usepackage{ctex}
\usepackage{amsmath,amssymb,amsfonts,mathtools}
\usepackage[mathbf=sym]{unicode-math}
\setmainfont{STIX Two Text}
\setmathfont{STIX Two Math}
\setsansfont{Arial}
\setmonofont{Courier New}
\setCJKmainfont[BoldFont=Noto Serif CJK SC Bold]{Noto Serif CJK SC}
\setCJKsansfont[BoldFont=Noto Sans CJK SC Bold]{Noto Sans CJK SC}
\setCJKmonofont{Noto Sans CJK SC}
\newCJKfontfamily\examkaishu{FandolKai-Regular}
\usepackage{graphicx}
\usepackage{booktabs,array,multirow,tabularx}
\usepackage{fancyhdr}
\usepackage{needspace}
\usepackage{tikz}
\usepackage{hyperref}
\hypersetup{colorlinks=false}
\usepackage{parskip}
\setlength{\parskip}{0.4em}
\setlength{\parindent}{0em}
\usepackage{xcolor}

% 数学符号
\renewcommand{\ge}{\geqslant}
\renewcommand{\geq}{\geqslant}
\renewcommand{\le}{\leqslant}
\renewcommand{\leq}{\leqslant}

% 知识点标签
\newcommand{\kp}[1]{%
    \par\vspace{0.3em}\noindent%
    {\small\textcolor{gray}{#1}}%
    \par\vspace{0.2em}%
}

% 答案解析区域
\newenvironment{solutionblock}{%
    \par\vspace{0.5em}%
    \noindent\rule{\textwidth}{0.4pt}%
    \vspace{0.5em}%
    \noindent\textbf{【参考答案与解析】}%
    \par\vspace{0.3em}%
    {\small%
}{%
    \par}%
}

% 答案标签
\newcommand{\anslabel}[1]{%
    \par\vspace{0.2em}\noindent%
    \textbf{#1}%
    \ignorespaces%
}

% 页眉页脚
\pagestyle{fancy}
\fancyhf{}
\fancyhead[L]{\small\examkaishu 虹口区高三二模·解析几何}
\fancyhead[R]{\small\examkaishu 第 \thepage 页}
\fancyfoot[C]{}
\renewcommand{\headrulewidth}{0.4pt}

\begin{document}

{\centering\Large\bfseries 四、解析几何\par}
\vspace{0.3em}
\hrule
\vspace{0.5em}

# 2025-2026学年虹口区高三二模数学试卷
## 四、解析几何
(椭圆综合)


\textbf{考试时间：120分钟  满分：150分}


20. 设椭圆 $\Gamma  : \frac{{x}^{2}}{4} + {y}^{2} = 1$ 的左顶点为 $A$ .

(1)求 $\Gamma$ 的离心率；

(2)设 $\Gamma$ 的左焦点为 $F$ ，上顶点为 $B$ ，若点 $P$ 在 $\Gamma$ 上且位于 $y$ 轴右侧. $\overrightarrow{AB}//\overrightarrow{FP}$ ，求点 $P$ 的横坐标;

(3)设直线 $l : x - y + m = 0$ ， $l$ 与 $\Gamma$ 交于不同的两点 $C$ 和 $D$ ，若点 $A$ 在以 ${CD}$ 为直径的圆外,求实数 $m$ 的取值范围.




<br><br><br><br><br><br><br><br>


---

\begin{solutionblock}
{\small \anslabel{【考点】：} 椭圆离心率、向量平行、直线与椭圆位置关系、圆}
{\small \anslabel{【答案】} (1) $\frac{\sqrt{3}}{2}$}
{\small (2) $x = \frac{\sqrt{5} - \sqrt{3}}{2}$}
{\small (3) $m \in  \left( {-\sqrt{5},\frac{6}{5}}\right)  \cup  \left( {2,\sqrt{5}}\right)$}
{\small \anslabel{【解析】}}
{\small \anslabel{【解析】}}
{\small \textbf{【小问 1 详解】}}
{\small 由椭圆方程可得 $a = 2, b = 1, c = \sqrt{{a}^{2} - {b}^{2}} = \sqrt{3}$ ,所以 $e = \frac{c}{a} = \frac{\sqrt{3}}{2}$ .}
{\small \textbf{【小问 2 详解】}}
{\small 由条件可知 $A\left( {-2,0}\right) , B\left( {0,1}\right) , F\left( {-\sqrt{3},0}\right)$ ,}
{\small 设直线 ${AB}$ 的斜率为 ${k}_{AB}$ ,直线 ${FP}$ 的斜率为 ${k}_{FP}$ ,}
{\small ${k}_{AB} = \frac{1 - 0}{0 + 2} = \frac{1}{2}$ ,因为 $\overrightarrow{AB}//\overrightarrow{FP}$ ,所以 ${k}_{FP} = \frac{1}{2}$ ,}
{\small 所以直线 ${FP}$ 的方程为 $y = \frac{1}{2}\left( {x + \sqrt{3}}\right)$ ,}
{\small 联立椭圆: $\left\{  {\begin{matrix} \frac{{x}^{2}}{4} + {y}^{2} = 1 \\  y = \frac{1}{2}\left( {x + \sqrt{3}}\right)  \end{matrix} \Rightarrow  {x}^{2} + {\left( x + \sqrt{3}\right) }^{2} = 4 \Rightarrow  2{x}^{2} + 2\sqrt{3}x - 1 = 0}\right.$ ,}
{\small 所以 $x = \frac{\sqrt{5} - \sqrt{3}}{2}$ 或 $x = \frac{-\sqrt{3} - \sqrt{5}}{2}$ ,}
{\small 又因为点 $P$ 位于 $y$ 轴右侧,所以 $P$ 的横坐标为 $x = \frac{\sqrt{5} - \sqrt{3}}{2}$ .}
{\small \textbf{【小问 3 详解】}}
{\small 设 $C\left( {{x}_{1},{y}_{1}}\right) , D\left( {{x}_{2},{y}_{2}}\right)$ ,联立椭圆}
{\small $\left\{  {\begin{array}{l} x - y + m = 0 \\  \frac{{x}^{2}}{4} + {y}^{2} = 1 \end{array} \Rightarrow  \frac{{\left( y - m\right) }^{2}}{4} + {y}^{2} = 1 \Rightarrow  5{y}^{2} - {2my} + {m}^{2} - 4 = 0,}\right.$}
{\small 先确定有两个交点,即 $\Delta  = {\left( -2m\right) }^{2} - 4 \times  5 \times  \left( {{m}^{2} - 4}\right)  > 0$ ,}
{\small 即 $4{m}^{2} - {20}{m}^{2} + {80} > 0 \Rightarrow   - {16}{m}^{2} + {80} > 0 \Rightarrow  {m}^{2} < 5$ ,}
{\small 所以 $- \sqrt{5} < m < \sqrt{5}$ .}
{\small 因为圆上任意一点与直径两端点连线所成的角为直角,}
{\small 而点 $A$ 在以 ${CD}$ 为直径的圆外,所以 $\angle {CAD} < {90}^{ \circ  }$ ,等价于 $\overrightarrow{AC} \cdot  \overrightarrow{AD} > 0$ ,}
{\small 由 $\overrightarrow{AC} = \left( {{x}_{1} + 2,{y}_{1}}\right) ,\overrightarrow{AD} = \left( {{x}_{2} + 2,{y}_{2}}\right)$ ,}
{\small 所以 $\overrightarrow{AC} \cdot  \overrightarrow{AD} = \left( {{x}_{1} + 2}\right) \left( {{x}_{2} + 2}\right)  + {y}_{1}{y}_{2} > 0$ ,}
{\small 即 $2{y}_{1}{y}_{2} - \left( {m - 2}\right) \left( {{y}_{1} + {y}_{2}}\right)  + {\left( m - 2\right) }^{2} > 0$ .}
{\small 将 ${y}_{1} + {y}_{2} = \frac{2m}{5},{y}_{1}{y}_{2} = \frac{{m}^{2} - 4}{5}$ 代入可得,}
{\small $2 \cdot  \frac{{m}^{2} - 4}{5} - \left( {m - 2}\right)  \cdot  \frac{2m}{5} + {\left( m - 2\right) }^{2} > 0 \Rightarrow  2\left( {{m}^{2} - 4}\right)  - {2m}\left( {m - 2}\right)  + 5{\left( m - 2\right) }^{2} > 0$ ,}
{\small $2{m}^{2} - 8 - 2{m}^{2} + {4m} + 5\left( {{m}^{2} - {4m} + 4}\right)  > 0 \Rightarrow  5{m}^{2} - {16m} + {12} > 0 \Rightarrow  \left( {{5m} - 6}\right) \left( {m - 2}\right)  > 0$ ,}
{\small 解得 $m < \frac{6}{5}$ 或 $m > 2$ ,结合 $- \sqrt{5} < m < \sqrt{5}$ ,}
{\small 所以 $m \in  \left( {-\sqrt{5},\frac{6}{5}}\right)  \cup  \left( {2,\sqrt{5}}\right)$ .}
{\small ![bo_d7obids91nqc738536r0_25_241_1031_388_283_0.jpg](images/bo_d7obids91nqc738536r0_25_241_1031_388_283_0.jpg)}
\end{solutionblock}
\end{document}